\chapter{System Overview}

The proposed system is a local web-based prototype for cardiopulmonary sound separation. A user uploads a mixed WAV audio file through the browser interface, and the FastAPI backend validates the upload, stores the original file, creates database records, runs NeoSSNet separation, saves the separated heart and lung audio files, and returns results for preview and download. SQLite is used for structured records such as uploaded audio metadata, model information, separation jobs, result paths, metrics, and logs, while local folders are used for the actual audio files and model checkpoints.

The following diagrams describe the system from different software design viewpoints. The use case diagram presents the system functions from the user's perspective, the class diagram shows the planned software components, the entity relationship diagram shows the database structure, and the object diagram gives an example of one completed separation job.

\section{Generic Use Case}

\begin{figure}[H]
  \centering
  \includegraphics[width=0.95\textwidth]{use_case_diagram.png}
  \caption{Generic Use Case Diagram}
  \label{fig:generic-use-case}
\end{figure}

Figure~\ref{fig:generic-use-case} shows the main actions available to a single generic actor, the User. The diagram is intentionally simple because the same user-facing workflow can be demonstrated by a student, lecturer, examiner, or researcher without changing the system boundary. The central use cases are uploading mixed audio, running cardiopulmonary sound separation, previewing separated heart and lung outputs, downloading separated outputs, and viewing processing history.

The simplified diagram keeps internal responsibilities such as validation, model selection, metadata storage, and logging out of the actor view. Those responsibilities are still required by the backend, but they support the visible workflow rather than appearing as separate user goals. This makes the use case model easier to read while still showing how upload, separation, preview, download, and history are connected.

\begin{landscape}
\begin{table}[H]
  \centering
  \scriptsize
  \setlength{\tabcolsep}{2pt}
  \caption{Use Case Descriptions}
  \label{tab:use-case-descriptions}
  \begin{tabular}{@{}p{0.11\linewidth} p{0.06\linewidth} p{0.15\linewidth} p{0.12\linewidth} p{0.15\linewidth} p{0.13\linewidth} p{0.13\linewidth}@{}}
    \toprule
    \textbf{Use Case Name} & \textbf{Actor} & \textbf{Description} & \textbf{Precondition} & \textbf{Main Flow} & \textbf{Alternative Flow} & \textbf{Postcondition} \\
    \midrule
    Upload Mixed Audio & User & Uploads a mixed cardiopulmonary WAV file through the browser interface. & The user has a supported WAV file. & The user selects a file and submits it; the backend validates and stores the upload. & If the file is invalid, the system rejects it and shows an error message. & Uploaded audio metadata and file path are stored. \\
    Run Sound Separation & User & Starts separation for an uploaded audio file using the selected or active model. & A valid uploaded audio record exists. & The user starts separation; the backend creates a job, runs the model, and stores output metadata. & If processing fails, the job is marked failed and the error is logged. & Heart and lung output files are generated for a successful job. \\
    Preview Separated Audio & User & Plays the separated heart and lung outputs in the browser. & A separation job has completed successfully. & The user opens the result view and plays the available audio controls. & If a file is missing or unsupported by the browser, the system keeps the download option available. & The user can inspect the separated outputs without leaving the interface. \\
    Download Separated Audio & User & Downloads the generated heart and lung WAV files. & A completed result contains valid output paths. & The user selects the heart or lung download link. & If the requested file is unavailable, the system returns a clear error response. & The selected separated audio file is saved by the user's browser. \\
    View Processing History & User & Reviews previous uploads, jobs, statuses, models, and result links. & At least one upload or job record exists, or the history list may be empty. & The user opens the history section and the system loads stored records from SQLite. & If no records exist, the system displays an empty history state. & The user can trace previous processing activity and reopen results. \\
    \bottomrule
  \end{tabular}
\end{table}
\end{landscape}

\section{Class Diagram}

\begin{landscape}
\begin{figure}[H]
  \centering
  \includegraphics[width=1.25\textwidth]{overall_class_diagram.png}
  \caption{Overall System Class Diagram}
  \label{fig:overall-class-diagram}
\end{figure}
\end{landscape}

Figure~\ref{fig:overall-class-diagram} presents the planned software classes grouped into frontend, router, service, machine learning, and persistence responsibilities. The browser interface calls FastAPI routers, the routers delegate work to service classes, and the service layer coordinates storage, model selection, algorithm creation, machine learning inference, and SQLite result records.

The class diagram also identifies where the selected design patterns are located:
\begin{itemize}
  \item \texttt{SeparationService} is the Facade because it presents one separation workflow to the router.
  \item \texttt{SeparationAlgorithm} and \texttt{NeoSSNetStrategy} form the Strategy structure because the service calls the algorithm through an abstraction.
  \item \texttt{SeparationAlgorithmFactory} is the Factory Method component because it creates the correct algorithm object from the selected model record.
\end{itemize}
This design keeps upload, separation, storage, result viewing, and history modular while avoiding hardcoded NeoSSNet logic in the router.

\section{Entity Relationship Diagram (ERD)}

\begin{figure}[H]
  \centering
  \includegraphics[width=0.82\textwidth]{erd.png}
  \caption{Entity Relationship Diagram}
  \label{fig:erd}
\end{figure}

Figure~\ref{fig:erd} shows the SQLite database structure used to store structured project data. The \texttt{uploaded\_audio} table records each uploaded WAV file, the \texttt{model} table records available model information, and the \texttt{separation\_job} table connects an uploaded file with the model used for processing. The \texttt{separation\_result} table stores the paths and metadata for generated heart and lung outputs.

The design intentionally stores file paths and metadata in SQLite rather than storing large WAV audio blobs in the database. This supports a practical student project architecture: local storage keeps uploaded and generated audio files, while SQLite supports history, result lookup, job status, metrics, and logs. As a result, the browser can display previous jobs and download links without scanning folders manually.

\section{Object Diagram}

\begin{figure}[H]
  \centering
  \includegraphics[width=0.88\textwidth]{object_diagram.png}
  \caption{Object Diagram for a Completed Separation Job}
  \label{fig:object-diagram}
\end{figure}

Figure~\ref{fig:object-diagram} gives a concrete example of how the system data may look after one successful separation. One \texttt{UploadedAudio} object is linked to one \texttt{SeparationJob}, the job uses an active \texttt{Model}, and the completed job produces a \texttt{SeparationResult} containing heart and lung output file paths.

The object diagram helps explain how upload, separation, storage, result viewing, and history are connected at runtime. When a user views history, the system can retrieve the job, original upload metadata, model information, output paths, optional metrics, and log record as related objects. This gives the project a traceable flow from the original mixed input to the separated outputs.
